\clearpage\begingroup
\NoSourceHead\SourceLeaf{183}
\begin{center}
{\Large\textls[150]{NACHLASS.}}\par\bigskip
ZURÜCKFÜHRUNG\par
DER WECHSELWIRKUNGEN ZWISCHEN GALVANISCHEN STRÖMEN\par
UND MAGNETISMUS AUF ABSOLUTE MAASSE.\par\bigskip
\textbf{[Erster Abschnitt.]}\par\medskip
1.
\end{center}
\noindent\hspace*{\parindent}Alle rein magnetischen Phänomene lassen sich durch die Annahme zweier magnetischen Fluida erklären, die nach einfachen der allgemeinen Schwere analogen Gesetzen auf einander wirken. Auch die Wechselwirkungen zwischen Magnetismus und Galvanismus nöthigen uns nicht, von jener Vorstellungsart abzugehen: allein wir müssen den früher bekannt gewesenen Naturgesetzen, nach welchen allein wir jene Phänomene noch nicht erklären können, andere beifügen, die man vor der Hand wenigstens wie neue physische Thatsachen zu betrachten hat, ohne die Hoffnung aufzugeben, dass es in Zukunft gelingen werde, sie auf noch einfachere Grundgesetze zurückzuführen.

\begin{center}2.\end{center}
\noindent\hspace*{\parindent}Wie es bei jeder schwierigen und verwickelten Naturwirkung nothwendig ist, mit den einfacheren Fällen anzufangen, so werden wir uns in der gegenwertigen Untersuchung auf die Betrachtung der galvanischen Ströme in solchen Leitern beschränken, deren Querschnitte sehr klein sind, so dass wir diese Querschnitte wie physische Punkte ansehen, oder die Ungleichheit der Entfernungen der einzelnen Punkte eines Querschnittes von einem äussern Punkte, auf welchen der Strom wirkt, bei Seite setzen dürfen. Der Stromträger kann
\clearpage\SourceLeaf{184}\HistoricalHead{179}{R}{WECHSELWIRKUNGEN ZWISCHEN GALVANISCHEN STRÖMEN UND MAGNETISMUS.}
\noindent daher wie eine mathematische Linie betrachtet werden, deren Länge, von einem willkürlichen Anfangspunkte bis zu jedem unbestimmten Punkte in dem Sinn des positiven Stroms gemessen, wir durch $s$ bezeichnen: die positive Richtung ist die bei hydrogalvanischer Erregung vom Zink durch das Wasser nach dem Kupfer gehende.

\begin{center}3.\end{center}
\noindent\hspace*{\parindent}Das Mittel zur Abmessung der Stärke eines galvanischen Stroms finden wir in dessen Wirkung auf die magnetischen Fluida. Das Grundgesetz dafür, welches bald nach \textsc{Oersteds} Entdeckung durch vielfache Versuche gefunden und bestätigt wurde, ist folgendes: Das Strom-Element in dem Theile des Leiters $AB=ds$ übt auf ein magnetisches Element $\mu$ in dem Punkte $C$ eine bewegende Kraft in einer Richtung aus, welche nicht in der geraden Linie $AC$ liegt, sondern gegen die durch $ABC$ gelegte Ebene normal ist, und zwar steht diese bewegende Kraft im zusammengesetzten geraden Verhältniss der Länge $ds$, der Intensität des Stromes in $A$, welche wir durch $g$ bezeichnen wollen, der Menge des freien magnetischen Fluidums in $\mu$, und des Sinus des Winkels $BAC=u$, endlich daneben im verkehrten doppelten Verhältniss der Entfernung $AC=r$, und kann also durch die Formel
\[
\frac{\mu g\sin u\mathbin{.}ds}{rr}
\]
\noindent ausgedrückt werden. Dabei ist folgendes zu bemerken:

I. Läge $C$ in der durch $AB$ gehenden geraden Linie selbst, so bliebe zwar die Ebene $ABC$ unbestimmt, allein in diesem Falle findet wegen $\sin u=0$ gar keine Kraft Statt.

II. In allen übrigen Fällen gehen von $C$ \textls[120]{zwei} Normallinien gegen das Planum $ABC$ aus, $CD$ und $CD'$, wovon die eine der Richtung der Kraft entspricht, welche das Stromelement auf ein Element positiven magnetischen Fluidums ausübt, während ein negatives magnetisches Molecül in der entgegengesetzten sollicitirt wird. Die Unterscheidung ist eine völlig bestimmte, wie die von der Ordnung unter je dreien von einem Punkte ausgehenden und nicht in Einer Ebene liegenden geraden Linien, welche jedoch nicht durch Begriffe a priori, sondern nur durch Hinweisung auf eine in der wirklichen Welt bestehende in der Anschauung gegeben werden kann. Der Erfahrung
\HistoricalFoot{}{23\textsuperscript{*}}
\clearpage\SourceLeaf{185}\HistoricalHead{180}{L}{PHYSIK.\quad NACHLASS.}
\noindent zufolge verhalten sich in dieser Beziehung die drei Richtungen von $A$ nach $B$, von $A$ nach $C$ und von $C$ nach $D$ (oder was dasselbe ist von $A$ parallel mit $CD$) wie diejenige, die wir durch die Benennung Richtung nach vorn, Richtung nach der rechten Seite, und Richtung nach unten unterscheiden, wenn $D$ sich auf positiven Magnetismus bezieht, während in derselben Ordnung die Richtung nach oben der auf den negativen Magnetismus wirkenden Kraft entspricht.

III. Da die Formel $\frac{\mu g\sin u\mathbin{.}ds}{rr}$ eine bewegende Kraft ausdrückt, eben so wie $\frac{\mu\mu}{rr}$, welches die bewegende Kraft ist, welche ein dem $\mu$ gleiches magnetisches Element auf $\mu$ ausübt, so folgt, dass $gds$ mit $\mu$ homogen ist, oder die Intensität $g$ homogen mit einer beliebigen Menge magnetischen Fluidums dividirt durch eine Linie.

Genau genommen sind alle diese Grössen Zahlen, welche ihr Verhältniss zu beliebigen Grössen derselben Art, welche man als Einheiten wählt, ausdrücken. Für die Intensität $g$ liegt dann als Einheit die Intensität eines solchen Stromes zum Grunde, welcher einen Kreis bildet, dessen Halbmesser die Längeneinheit ist, und von welchem ein dem Halbmesser gleicher Bogen in seinem Mittelpunkt auf magnetische Fluida eben so stark wirkt, wie die als Einheit angenommene Menge des magnetischen Fluidums in derselben Entfernung wirken würde. Es wird nicht überflüssig sein, hier diejenigen Formeln zusammenzustellen, welche die Homogeneität der verschiedenen bei diesen und ähnlichen Untersuchungen vorkommenden Grössen darstellen. Es sei $R$ eine Lineargrösse, $t$ eine Zeitgrösse, $p$ eine Masse, $\mu$ ein Quantum magnetischen Fluidums, $T$ die Kraft des Erdmagnetismus oder eine damit homogene Kraft. Es ist sodann homogen
\[
\begin{aligned}
\mu&\text{ mit }\surd\frac{pR^3}{tt}\\[5pt]
T&\text{ mit }\surd\frac{p}{Rtt}\\[5pt]
g&\text{ mit }\frac{\mu}{R}\text{ oder mit }\surd\frac{pR}{tt}\text{ oder mit }RT.
\end{aligned}
\]
\noindent Homogene Grössen werden immer durch Zahlen ausgedrückt, deren Verhältniss von der Wahl der Fundamentaleinheiten für Raum, Zeit und Masse unabhängig ist.

IV. Nach allen Erfahrungen ist die Intensität eines galvanischen Stroms
\clearpage\SourceLeaf{186}\HistoricalHead{181}{R}{WECHSELWIRKUNGEN ZWISCHEN GALVANISCHEN STRÖMEN UND MAGNETISMUS.}
\noindent an allen Stellen dieselbe, insofern keine Theilung Statt findet. Bei einer Theilung in zwei oder mehrere Äste ist die Summe der Intensitäten in den einzelnen Ästen der Intensität vor der Theilung gleich; etwas ähnliches findet Statt bei der Vereinigung mehrerer Äste zu Einem Strome, dessen Intensität der Summe der Intensitäten in den einzelnen vereinigten Armen gleich wird.

V. Zwei Ströme von gleicher Intensität, die eine und dieselbe Linie in entgegengesetzter Richtung durchlaufen, bringen in jedem Punkte des Raumes gleiche aber entgegengesetzte bewegende Kräfte hervor, mithin gar keine, wenn jene als coexistirend gedacht werden.

VI. Gegenseitigkeit oder Rückwirkung auf das Element des Stromträgers[*)].\begingroup\renewcommand{\thefootnote}{}\footnotetext{[*) Diese Zeile steht am Rande der Handschrift; die weitere Ausführung fehlt.]}\endgroup

\begin{center}4.\end{center}
\noindent\hspace*{\parindent}Nach dem im vorigen Artikel aufgestellten Grundgesetz wird nun die Bestimmung der Kraft, welche ein Strom in jedem Punkte des Raumes auf die magnetischen Fluida ausübt, zu einer rein mathematischen Aufgabe. Ehe wir jedoch dafür allgemeine Formeln entwickeln, wollen wir einige der einfachsten Fälle besonders betrachten. Wir werden dabei den Factor $g\mu$ weglassen oder die Wirkung des Stroms, dessen Intensität $=1$ ist, auf die in Einem Punkte vereinigt gedachte Einheit des positiven magnetischen Fluidums zum Grunde legen.

Wir betrachten zuerst einen geradlinichten Strom, der auf einen Punkt $P$ wirkt; es sei $PA=a$ senkrecht gegen die gerade Linie und $x$ ein unbestimmtes Stück der letzteren, von $A$ an in dem Sinne der Richtung des Stromes gezählt. Der Winkel, den jedes Stromelement $dx$ mit der von seinem Anfangspunkt nach $P$ gezogenen geraden Linie $r$ macht, sei $u$, also
\[
\frac{\sin u\mathbin{.}dx}{rr}
\]
\noindent die Wirkung des Stromelements $dx$ auf $P$, welche dem vorhergehenden zufolge senkrecht gegen die Ebene durch $P$ und die gerade Linie, und zwar nach unten zu gerichtet ist, wenn man sich diese Ebene horizontal denkt und $P$ rechts von der geraden Linie liegt, die Richtung des Stromes in dieser als vorwärtsgehend betrachtet.

\clearpage\SourceLeaf{187}\HistoricalHead{182}{L}{PHYSIK.\quad NACHLASS.}
Da
\[
x=-a\operatorname{cotg}u,\quad r=+\frac{a}{\sin u},
\]
\noindent so wird
\[
\int\frac{\sin u\,dx}{rr}=\int\frac{\sin u\,du}{a}=\text{Const.}-\frac{\cos u}{a}.
\]
\noindent Sind also die Werthe von $x,u,r$ in dem Anfangs- und Endpunkte des Stromes resp.\ $x^0,u^0,r^0;\ x',u',r'$, so wird die ganze in $P$ wirkende Kraft
\[
=\frac{\cos u^0-\cos u'}{a}=\frac{x'}{ar'}-\frac{x^0}{ar^0}.
\]
\noindent Erstreckt sich die gerade Linie auf einer oder beiden Seiten ins unendliche, so bleibt demungeachtet die Kraft endlich; es wird dieselbe nemlich
\[
\begin{aligned}
&=\frac{1-\cos u'}{a}=\frac{r'+x'}{ar'}\text{[*)], wenn der Strom ohne Anfang,}\\[5pt]
&=\frac{1+\cos u^0}{a}=\frac{r^0-x^0}{ar^0}\text{[*)], wenn er ohne Ende,}\\[5pt]
&=\frac{2}{a},\text{ wenn er ohne Anfang und Ende ist,}
\end{aligned}
\]
\begingroup\renewcommand{\thefootnote}{}\footnotetext{[*) In der Handschrift befinden sich an dieser Stelle Schreibfehler.]}\endgroup
\noindent und es erhellet hieraus, dass alle über eine gewisse Entfernung hinaus liegenden Theile eines Stromes in einem Leitungsdraht keine merkliche Wirkung mehr ausüben.
\clearpage\endgroup
